\chapter{“बीस सेकंड”}
\label{sec:1}
एक इंटेलिजेंस अफ़सर ने, मशीन से बनी हत्या-सूचियों के इर्द-गिर्द खड़ी निशाना तय करने वाली एक पाइपलाइन में अपनी भूमिका बताते हुए, एक पत्रकार से कहा कि वह हर नाम पर बीस सेकंड लगाता था, बस इतना कि जाँच ले कि निशाना मर्द है: “किसी की आवाज़ मर्दाना है या ज़नाना, यह बताने में ज़्यादा वक़्त नहीं लगता।” उसने आगे कहा कि एक इंसान के तौर पर उसका योगदान शून्य था, सिवाय इसके कि वह मंज़ूरी का ठप्पा था। “इससे बहुत वक़्त बचता था” \autocite{abraham2024lavender}। उसके बयान के मुताबिक़, ऑपरेशन की रफ़्तार में “लूप में इंसान” का कुल मतलब यही था।\footnote{ऑटोमेशन पूर्वाग्रह पर हुआ शोध इसी तरह काम सौंप देने की भविष्यवाणी करता है: जब दूसरी माँगें लोगों का ध्यान खींच रही हों, तो जिन्हें मशीन की सिफ़ारिश दी जाती है, वे उसे कम जाँचते हैं और उसकी बात ज़्यादा मान लेते हैं \autocite{parasuraman2010complacency}।}

एक व्याख्या यह है कि अफ़सर एक ऐसा कर्ता है जिसे बहुत कम समय दिया गया, और इलाज यह है कि उसे बेहतर ढंग से चलाया जाए — ज़्यादा समय, एक दूसरा समीक्षक, लूप में बिठाया गया कुछ और ज़मीर। इस \emph{सामाजिक अधीनता} वाले नज़रिए से इंसान एक कर्ता है जो मशीन के बाहर खड़ा होकर उसे औज़ार की तरह चलाता है: ड्राइवर और कार, मज़दूर और खराद। इंसान की एक भूमिका है, एक नाम है, ज़िम्मेदारियाँ हैं, और एक ज़मीर है जिससे बात की जा सकती है। दूसरा नज़रिया \emph{मशीनी ग़ुलामी} का है, जिसमें इंसान मशीन का एक पुर्ज़ा है, जिसे मशीन के दूसरे हिस्सों के साथ जोड़कर एक ऐसा सिस्टम बनाया गया है जो उन सबको एक साथ नियंत्रित करता है; मशीन इंसान को अपने एक अंग की तरह इस्तेमाल करती है, और इंसान का योगदान एक काम है जो वह सिस्टम की तय की हुई रफ़्तार पर करता है।\footnote{डॉक्टरो अफ़सर की इस स्थिति को उल्टा सेंटॉर कहते हैं: ऐसा इंसान नहीं जिसकी मदद मशीन करती है, बल्कि इंसान से बना मशीन का एक पुर्ज़ा \autocite{doctorow2026reversecentaur}।} ये दोनों एक-दूसरे को काटते नहीं; कोई इंसान एक स्तर पर भूमिका निभाते हुए दूसरे स्तर पर एक हिस्सा हो सकता है \autocite[पठार 13, “The Apparatus of Capture”: ग़ुलामी तब है जब इंसान ख़ुद एक ऐसी मशीन के घटक पुर्ज़े होते हैं जिसे वे आपस में और दूसरी चीज़ों के साथ मिलकर बनाते हैं, किसी ऊँची इकाई के नियंत्रण और निर्देशन में; और अधीनता तब है जब वह ऊँची इकाई इंसान को एक कर्ता के रूप में गढ़ती है जो किसी नई बाहरी वस्तु से जुड़ा है, और वह वस्तु कोई जानवर, कोई औज़ार, या यहाँ तक कि कोई मशीन भी हो सकती है]{deleuze1987plateaus}, और बीस सेकंड वाला अफ़सर दोनों था। एक हिस्से के रूप में, वह वह जगह था जहाँ पाइपलाइन अपना ठप्पा पैदा करती थी, और बीस सेकंड उस हिस्से का चक्र-समय था। एक कर्ता के रूप में, उससे पूरे समय बात की जाती रही: उसने एक पत्रकार को बताया कि एक तरफ़ आम नागरिकों की मौतों की ख़बरें थीं और दूसरी तरफ़ काम पर और निशानों का दबाव, और इन दोनों के बीच उसके भीतर कैसा टकराव था। और यह टकराव उसके ज़मीर तक पहुँचा, फिर भी वह साथ चलता रहा। टकराव दर्ज हुआ और उसने कुछ तय नहीं किया;\footnote{या कम से कम ऐसा कुछ नहीं जिसने इज़राइली रक्षा बलों को अक्टूबर 2023 से अक्टूबर 2025 के बीच ग़ज़ा में एक लाख लोगों को मारने से रोका हो \autocite{gomezugarte2025uncertainty,acosta2026gaza,economist2025gaza}। और उस ख़ास अफ़सर का बयान अप्रैल 2024 में +972 Magazine ने छापा था। कोई अंदाज़ा ही लगा सकता है कि उस अफ़सर के एक पत्रकार से बात करने के फ़ैसले ने एक और मौत रोकी, 50,000 और मौतें रोकीं, या एक भी नहीं। और अगर अफ़सर उस भूमिका में दो हफ़्ते बाद इनकार कर देता तो क्या होता — मान लीजिए, नवंबर 2023 में? हम एक ऐसे इंटेलिजेंस अफ़सर की बात कर रहे हैं जो Lavender की मदद से हत्या-सूचियाँ तैयार करता था। उसका रोज़ का काम का माहौल शायद किसी एक्शन फ़िल्म से ज़्यादा किसी आम दफ़्तरी नौकरी जैसा था। क्या उसे कमान की सीढ़ी में ऊपर तक अपने अफ़सरों से बार-बार मिलने की कोशिश करनी चाहिए थी, और साथ ही ऐसे हवाई हमलों के निशानों को मंज़ूरी देने से इनकार करना चाहिए था जिनसे आम नागरिकों की मौत की उचित आशंका थी? कहा जा सकता है कि यही नैतिक रास्ता होता। उस काल्पनिक स्थिति में, क्या कम लोग मरते?} यह फ़ैसला कि नीति अनुपात से बाहर थी, उसके मुताबिक़ बाद में, पीछे मुड़कर देखने पर आया।

तो जो चीज़ ग़ायब थी, वह ज़मीर नहीं था; ज़मीर तो उसके पास था। ज़मीर किसी फ़ैसले की कुछ क़ीमत तय करके काम करता है, पर यहाँ क़ीमतें उल्टी तरफ़ बैठती थीं: “हम पर लगातार दबाव डाला जाता था: ‘और निशाने लाओ।’ वे सचमुच हम पर चिल्लाते थे” \autocite{abraham2024lavender}। वह संकोच वाला एक पुर्ज़ा था, पर रफ़्तार बनी रही: उसके संकोच से उसके योगदान में कुछ नहीं बदला, और वह एक ठप्पा बन गया। संकोच वाला जो पुर्ज़ा रफ़्तार नहीं थाम पाता, वह वही “इंसानी अड़चन” बन जाता है जिसे हटाने पर इज़राइल की सिग्नल-इंटेलिजेंस यूनिट के एक पूर्व कमांडर ने किताब लिखी \autocite{ys2021humanmachine}। अफ़सर को पुर्ज़ा उसके रवैये ने नहीं, इंतज़ाम ने बनाया था। पुर्ज़ा अनमने ढंग से भी साथ चल सकता है और उत्साह से भी: आरेंट ने आइखमान के मुक़दमे में जिस बिना सोचे-समझे काम करने वाले कारिंदे को देखा था \autocite{arendt1963eichmann}, वह, उन बातों के सबूत पर जो उसने चार साल पहले अर्जेंटीना में बिना किसी दाँव के कही थीं, यरूशलम के लिए किया गया एक दिखावा था \autocite{stangneth2014eichmann}। पाइपलाइन ने अफ़सर को जो योगदान देने दिया, वह एक तय रफ़्तार पर एक तय काम था: उसने हर मामले को वैसे ही देखा जैसे पाइपलाइन ने उसे बताया, एक बार में एक, उतने समय में जितना पाइपलाइन ने उसे दिया। क्या कोई चीज़ उस जगह खड़ी होकर भी पुर्ज़ा न बने, यह भाग IV का सवाल है।

खोजी पत्रकारिता ने ग़ज़ा में साथ-साथ काम करने वाले तीन सिस्टमों का ब्योरा दिया: पहला इमारतों और बुनियादी ढाँचे के निशाने बनाता था, और सूत्रों के मुताबिक़ उसके ज़्यादातर सुझाव निजी घर थे; दूसरा लोगों को उनके अनुमानित उग्रवादी जुड़ाव के आधार पर अंक देता था, और ऐसी सूचियाँ बनाता था जिनमें कथित तौर पर 37,000 तक लोगों को चिह्नित किया गया था, और भीतर यह माना जाता था कि वे दस में से लगभग एक बार ग़लत होती हैं\footnote{इज़राइली सेना ने इस ब्योरे को चुनौती दी। नीचे मीनाब और Maven के ब्योरे भी ऐसी खोजी पत्रकारिता पर टिके हैं जिसे संबंधित पक्ष चुनौती देते हैं। मैं उन्हें वैसे ही लेता हूँ जैसे वे रिपोर्ट किए गए।}; और तीसरा इस पर नज़र रखता था कि चिह्नित लोग कब घर लौटते हैं, ताकि हवाई हमले ठीक उस वक़्त किए जा सकें जब वे अपने परिवारों के साथ हों।

\runin{Maven: अमेरिकी संस्करण, अब जनरेटिव एआई की रफ़्तार के साथ}

मीनाब के स्कूल पर हुई बमबारी में 156 लोग मारे गए, जिनमें 120 स्कूली बच्चे थे \autocite[पैरा~92]{UNFFMIran2026}।\footnote{गिनती स्रोत और तारीख़ के हिसाब से बदलती है, 156 (मीनाब के अभियोजक की अंतिम गिनती, जो Amnesty International ने भी दी) से लेकर 168 तक (शहर के मेयर की, शुरुआती दिनों में)। संयुक्त राष्ट्र मिशन अभियोजक के आँकड़े देता है और साथ ही स्वतंत्र जानकारी भी बताता है कि 157 लोग मारे गए, जिनमें 123 बच्चे 13 साल या उससे कम उम्र के थे \autocite[पैरा~92–93]{UNFFMIran2026}।} स्कूल का नागरिक इस्तेमाल खुले स्रोतों वाली सैटेलाइट तस्वीरों में दिखता था, फिर भी निशाने की कोडिंग में अब भी बगल के रिवोल्यूशनरी गार्ड ठिकाने से विरासत में मिला पदनाम चला आ रहा था। संयुक्त राष्ट्र के एक तथ्य-खोजी मिशन ने कहा कि यह मानने के ठोस आधार हैं कि यह बमबारी एक नागरिक वस्तु को हमले का निशाना बनाने का युद्ध अपराध थी \autocite{ohchr2026minabpress}, और यह कि स्कूल साफ़ तौर पर पहचाना जा सकता था और ऐसी कोई जानकारी नहीं थी जो उसके सैन्य इस्तेमाल की ओर इशारा करे \autocite[पैरा~81–86, 124–126]{UNFFMIran2026}। मिशन ने पाया कि संयुक्त राज्य अमेरिका निशाने की पुष्टि के लिए हर संभव कोशिश करने के अपने दायित्व में नाकाम रहा, उसने निशाने के आँकड़ों को अद्यतन करने को संभव सावधानियों में गिना, और इस नाकामी को लापरवाही से आगे की चीज़ माना \autocite[पैरा~124–125]{UNFFMIran2026}। Amnesty International भी इसी तरह के नतीजे पर पहुँचा, ऐसी भाषा में जो अतिरिक्त प्रोटोकॉल I के अनुच्छेद 57(2)(a)(i) में दिए पुष्टि करने के कर्तव्य से मेल खाती है: कि सेनाओं ने शायद पुरानी ख़ुफ़िया जानकारी पर भरोसा किया और निशाने की पुष्टि के लिए हर संभव कोशिश नहीं की \autocite{Amnesty2026Minab}।

यह 28 फ़रवरी की सुबह हुआ, ईरान अभियान के पहले दिन \autocite{Amnesty2026Minab,crs2026iranwar}। उससे एक दिन पहले मंत्री ने ऐलान किया था कि जिस विक्रेता (Anthropic) का मॉडल (Claude) निशाना तय करने वाले प्लेटफ़ॉर्म (Maven Smart System) के भीतर बैठा था, उसे सप्लाई-चेन जोखिम घोषित किया जाएगा, क्योंकि उसने Claude को हर क़ानूनी इस्तेमाल के लिए लाइसेंस देने से इनकार किया था \autocite{npr2026openai}। Maven Smart System ने स्कूल को पहले दिन के लिए सुझाए गए निशाने के रूप में लौटाया था \autocite{livingstone2026maven,gizmodo2026minab}। कोई भी प्रकाशित ब्योरा यह नहीं बताता कि उस सिफ़ारिश में Claude की क्या भूमिका थी, और विक्रेता के काम से वाक़िफ़ एक व्यक्ति ने NBC News को बताया कि मॉडल ने सीधे निशाने तय करने की सलाह नहीं दी \autocite{nbc2026targeting}।

पेंटागन की एक अंदरूनी समीक्षा ने, जिसके पूरे निष्कर्ष अब भी सख़्ती से छिपाकर रखे गए हैं, कथित तौर पर पाया कि पिछले युद्धों में निशाने की जो तैयारी घंटों लेती थी, वह मिनटों में सिमट गई, कि नागरिक नुक़सान कम करने वाला स्टाफ़ लगभग नब्बे प्रतिशत घटकर बीस से भी कम लोगों पर आ गया, कि सेंट्रल कमांड की टीम दस से घटकर एक रह गई, और कि उसके किसी भी सदस्य ने मीनाब वाली जगह की समीक्षा नहीं की — एक नतीजा जिसके लिए अधिकारियों ने कटौतियों को भी ज़िम्मेदार ठहराया और इस फ़ैसले को भी कि टीम को योजना और ऑपरेशनों में शामिल न किया जाए \autocite{Bloomberg2026KillChain}। औपचारिक जाँच मार्च में शुरू हुई थी, जब एक शुरुआती आकलन ने पाया कि ज़िम्मेदार शायद संयुक्त राज्य अमेरिका है \autocite{NPR2026MinabProbe}।

2026 के ईरान अभियान में, ठेकेदार के बनाए एक निशाना-प्लेटफ़ॉर्म ने, जिसके भीतर एक व्यावसायिक विक्रेता का सामान्य-उपयोग वाला मॉडल था और जिसका इस्तेमाल निशानों को रणनीतिक अहमियत के हिसाब से क्रम देने और हर हमले के लिए क़ानूनी औचित्य का मसौदा लिखने में होता था \autocite{livingstone2026maven}, पहले चौबीस घंटों में एक हज़ार से ज़्यादा निशानों पर हमला करने में मदद की, जो उससे पहले संभव संख्या का लगभग दस गुना था \autocite{csis2026maven}, और अड़तीसवें दिन तक तेरह हज़ार पर \autocite{breakingdefense2026maven}। पेंटागन ने इसके लिए प्रति-निशाना समीक्षा का कोई समय प्रकाशित नहीं किया है।\footnote{Maven एक कार्यक्रम का नाम है जो एक से ज़्यादा सिस्टमों पर लागू होता है: पेंटागन का वह कार्यक्रम जिसे अब भी Maven कहा जाता है, और Maven Smart System, वह सॉफ़्टवेयर प्लेटफ़ॉर्म जो उसी से निकला \autocite{csis2026maven}। कार्यक्रम 2017 में एक परियोजना के रूप में शुरू हुआ था, जिसका काम इस्लामिक स्टेट के ख़िलाफ़ अभियान के दौरान ड्रोन जो वीडियो इकट्ठा कर रहे थे, उनमें चीज़ों पर लेबल लगाना था \autocite{weisgerber2017maven,fang2019maven}। 2026 का सिस्टम सैन्य सिस्टमों और ख़ुफ़िया स्रोतों को एक ही इंटरफ़ेस में जोड़ता है, जिससे निशाने ढूँढे और उन पर हमले किए जाते हैं और ऑपरेशनों की कमान सँभाली जाती है \autocite{defensescoop2026transition}। प्लेटफ़ॉर्म Palantir का है: पेंटागन के मुख्य डिजिटल और एआई अधिकारी इसे “Palantir का सॉफ़्टवेयर-ऐज़-ए-सर्विस प्रोडक्ट” कहते हैं \autocite{defensescoop2026transition}। Palantir का सार्वजनिक रुख़, जो दूसरे हाथ से उद्धृत है, यह है कि ज़िंदगी या मौत का फ़ैसला करना उसका काम नहीं और ज़िम्मेदारी हमेशा सैन्य संगठन की रहती है \autocite{lee2026palantir}, और वह कहता है कि बुनियादी ख़ुफ़िया जानकारी सरकार की ज़िम्मेदारी थी और इसका कोई सबूत नहीं कि Maven में कोई ख़राबी आई या उसने मीनाब की बमबारी करवाई \autocite{Bloomberg2026KillChain}। कार्यक्रम एक दफ़्तर और ठेकेदार से अगले तक कैसे पहुँचा, यह परिशिष्ट~\ref{sec:C.1} में है।}

इज़राइली और अमेरिकी हवाई हमलों की सूचियों में जो लोग थे, वे उन सूचियों में होने पर आपत्ति नहीं कर सकते थे — पाइपलाइनें इस तरह की प्रतिक्रिया लेने के लिए बनी ही नहीं थीं। आंद्री होसे एर्नांदेस रोमेरो एक अलग पाइपलाइन में फँसा था, और वह आपत्ति कर सकता था। वेनेज़ुएला का एक मेकअप आर्टिस्ट, वह अगस्त 2024 में शरण माँगने के लिए सैन इसिद्रो प्रवेश-द्वार पर ख़ुद हाज़िर हुआ, उसने अपना विश्वसनीय-भय इंटरव्यू पास किया, और उसके पास एक वकील, एक लंबित दावा और सुनवाई की एक तारीख़ थी \autocite{LULAC2025HernandezRomero,Vega2025HernandezRomero}। उसकी जाँच करने वाले एक ठेकेदार के कर्मचारी ने उसके दो टैटू, “माँ” और “पापा” के शब्दों के ऊपर बने मुकुटों को, गिरोह के निशान समझ लिया \autocite{Palomera2025Blade}, और मार्च 2025 में, उसकी सुनवाई से पहले, उसे विमान से अल साल्वाडोर की एक जेल में भेज दिया गया, उसी शाम जब एक संघीय जज ने विमानों को वापस मोड़ने का आदेश दिया था \autocite{jgg2025contempt}। आपत्ति करने के लिए जो कुछ चाहिए, वह सब उसके पास था, सिवाय समय के। जब हम संयुक्त राज्य अमेरिका की अपहरण और जबरन स्थानांतरण वाली पाइपलाइन पर बात करेंगे, तब उसके मामले पर विस्तार से लौटेंगे।

\textbf{किसने इनकार किया।} बीस सेकंड के रिकॉर्ड का एक दूसरा आधा हिस्सा है, उन लोगों का रिकॉर्ड जो इस तरह के दबाव के सामने टिके रहे, और वह बताता है कि कमरे की बनावट बहुत मायने रखती है। मिलग्राम की बुनियादी स्थिति में लगभग एक-तिहाई लोग सबसे बड़े झटके से पहले ही रुक गए। जब दो साथी अवज्ञा कर रहे थे, तब दस में से सिर्फ़ एक आख़िर तक गया, और जब दो प्रयोगकर्ता आपस में असहमत थे, तब बुनियादी स्थिति के मुक़ाबले कहीं कम लोग आगे बढ़े \autocite{milgram1974obedience}। ऐश के लगभग एक-चौथाई लोग ग़लत बहुमत के आगे कभी नहीं झुके \autocite{asch1956studies}। ब्राउनिंग का अनुमान है कि रिज़र्व पुलिस बटालियन 101 के दसवें से पाँचवें हिस्से तक लोगों ने यूज़ेफ़ुफ़ में, बटालियन के पहले नरसंहार में, गोली चलाने से इनकार किया या उससे बच निकले \autocite{browning1992ordinary}। माहौल के हिसाब से यह हिस्सा लगभग किसी के नहीं से लेकर लगभग सबके तक गया। जो लोग टिके रहे, उनके पास उस चैनल के बाहर से आया सबूत था: अवज्ञा करता एक साथी, असहमत होती एक दूसरी सत्ता, और, जिन बचाने वालों का अध्ययन ओलाइनर दंपति ने किया उनके बीच, ऐसे रिश्ते जो उस रेखा के आर-पार जाते थे जो उत्पीड़न ने खींची थी \autocite{oliner1988altruistic}। या उन्होंने अपनी प्रतिबद्धता को इतनी बारीकी से थामा कि दबाव उसे तोड़ न सका, जैसे नैतिक आदेश और पवित्र मूल्य थामे जाते हैं \autocite{skitka2009limits,ginges2007sacred}। या उन्होंने स्रोत को ऐसा माना जिसका अपना दाँव हो। और समूहों ने वहाँ विरोध किया जहाँ अकेले लोग शायद ही कभी कर पाए \autocite{gamson1982encounters}।

कमरे को दूसरे पक्ष के लिए भी उतनी ही आसानी से रचा जा सकता है। वाइमार जर्मनी में संगठनों और संस्थाओं का घना सामाजिक जीवन था, जिसने शायद गणराज्य को गिराने में मदद की: जैसे-जैसे पार्टियाँ और सरकारें लोगों की शिकायतों का जवाब देने में नाकाम हुईं, लोगों ने अपनी ऊर्जा समूहों की तर्ज़ पर बने क्लबों और संस्थाओं में लगा दी, और नाज़ियों ने उन्हीं नेटवर्कों से भर्ती की जो इन संस्थाओं ने पहले से बना रखे थे \autocite{berman1997weimar}। जिन क़स्बों में ज़्यादा गायक-मंडलियाँ, खेल क्लब और जिम्नास्टिक संघ थे, वहाँ लोग पार्टी में तेज़ी से शामिल हुए, पर सिर्फ़ वहाँ जहाँ सरकार अस्थिर थी; प्रशा में, जहाँ सरकार स्थिर थी, यह संबंध छोटा और नगण्य था \autocite{satyanath2017bowling}। तो बचाने वालों के पास जो था, वह सिर्फ़ ज़्यादा संगत नहीं थी। उनके रिश्ते उस रेखा के आर-पार जाते थे जो उत्पीड़न ने खींची थी। आशुतोष वार्ष्णेय ने भारतीय शहरों में यही पाया: हिंदुओं और मुसलमानों के बीच दंगे वहाँ कम होते थे जहाँ संगठन दोनों समुदायों को एक ही संस्थाओं में साथ लाते थे, और वहाँ आम थे जहाँ ऐसा नहीं था \autocite{varshney2002ethnic}। जो रिश्ते इनकार को संभव बनाते हैं, वही रिश्ते हैं जिन्हें किसी फ़ासीवादी आंदोलन को काटना पड़ेगा।

आगे जो कुछ है, वह कमरे को रचने के बारे में है।\footnote{कमरा वह सेट है जहाँ दृश्य होता है, और जैसा मेरे थिएटर के टीचर कहते, सेट ही दृश्य को निर्देशित करता है। कमरे में और कौन है? क्या वहाँ से निकलकर बचा जा सकता है? क्या कोई इनकार को सुनेगा?} रचना बनाने का ही एक शब्द है, और हन्ना आरेंट ने राजनीति को बनाई जाने वाली चीज़ मानने के ख़िलाफ़ चेताया था। उन्होंने लिखा कि यूनानियों के लिए क़ानून सचमुच एक दीवार था, जिसके बिना एक बस्ती तो हो सकती थी पर नगर नहीं \autocite{arendt1958human}: क़ानून बनाने वाले ने दीवार खड़ी की, और उसके भीतर जो होता था, वह उसके बनाने की चीज़ नहीं थी। जिन बुनियादी उपायों की मैं वकालत करता हूँ, वे इसी अर्थ में दीवारें हैं। वे इनकार को ऐसा बनाते हैं कि उसके बाद भी बचा जा सके, और उसे सुना जाए; वे उसे पैदा नहीं करते, और सत्ता, उनके हिसाब से, सिर्फ़ साथ मिलकर काम करते लोगों के बीच होती है और उनके बिखरते ही ख़त्म हो जाती है, और ‘किसने इनकार किया’ वाला रिकॉर्ड इनकार के बारे में यही दिखाता है। भाग II और III के औज़ार इनकार को गिनती में लाते हैं। भाग V की संस्थाएँ इनकार करने वाले बनाती हैं। इनकार कर सकने वाली मशीन को वही चाहिए होगा जो उन दस में से नौ के पास था जिनके बगल में एक साथी था: अपने बगल में एक अलग इतिहास वाला साथी, जाने के लिए कोई जगह, और कोई जो उसे सुने।

पाइपलाइन में इनकार करने वाला — इनकार करने वाली मशीन भी — अफ़सर के लिए यही हो सकता है। इनकार करने वाली मशीन तब स्क्रीन पर होती जब मामला उस तक पहुँचता: अवज्ञा करता साथी, अगर यह असर कमरे से स्क्रीन तक के सफ़र में बचा रहे, जिसे अभी परखा नहीं गया और जिसके बारे में ऑटोमेशन पूर्वाग्रह कहता है कि वह नहीं बचेगा। जो मशीन समीक्षकों को इस तरह मोड़ती कि वे उसकी बात मान लें, वह उन्हें बस दूसरी तरफ़ के ठप्पे बना देती। क्या बीस सेकंड और बगल में इनकार करती एक मशीन वाला अफ़सर वही अफ़सर है जो रिपोर्टिंग में है?
